\documentclass[12pt]{article}
\usepackage[utf8]{inputenc}
\usepackage[T1]{fontenc}
\usepackage{hyperref}
\title{Lingua Franca Language Resources for K-12 Curriculum Development}
\author{Compiled by Claude Sonnet 5.5 (claude-sonnet-5-5) for Andr\'as Kornai}
\date{}
\begin{document}
\maketitle

\section{Introduction}
Lingua Franca (Mediterranean Lingua Franca, also called Sabir; ISO 639-3 pml; Glottocode ling1242) is an extinct Romance-based pidgin used around the Mediterranean. The sources describe it as documented from about 1300 to the late 19th century, mostly in the Maghreb. The surviving corpus is small. No school materials, primers or revitalization programmes were found. This dossier lists scholarly sources only.

\section{Primary Educational Resources}
None found. There is no primer or textbook. Teaching would have to use the attested texts discussed in the scholarly works below.

\section{Linguistic and Reference Materials}
\begin{itemize}
\item Wikipedia, Mediterranean Lingua Franca (overview): \url{https://en.wikipedia.org/wiki/Mediterranean_Lingua_Franca}
\item Rizza et al., review article (PDF, University of Verona repository): \url{https://iris.univr.it/bitstream/11562/1153589/1/34_1-2_REZ_Rizza-et_al.pdf}
\item Language Science Press volume (catalogue entry): \url{https://langsci-press.org/catalog/view/235/2544/1852-1}
\item Operstein, article in the Italian Journal of Linguistics (PDF): \url{https://www.italian-journal-linguistics.com/app/uploads/2021/05/4_Operstein.pdf}
\item Baglioni (2024), ``What Remains of An Atypical Restsprache'' (PDF): \url{https://iris.unive.it/retrieve/d881d7be-1c05-4dd0-afc0-f2e1104fd937/Baglioni\%20\%282024\%29_What\%20Remains\%20of\%20An\%20Atypical\%20Restsprache.pdf}
\item Nolan, SOAS Working Papers in Linguistics vol. 17 (PDF): \url{https://eprints.soas.ac.uk/40403/1/SOAS\%20WPL\%20Vol\%2017\%20-\%20Nolan.pdf}
\item ``Lingua Franca of the Levant'' (Brill, PDF at Ca' Foscari repository): \url{https://iris.unive.it/retrieve/2951f9f4-e01e-4e9b-9816-a30abca468b9/Lingua\%20Franca\%20of\%20the\%20Levant\%20\%e2\%80\%94\%20Brill.pdf}
\item Cambridge University Press book excerpt (PDF): \url{https://assets.cambridge.org/97813165/18311/excerpt/9781316518311_excerpt.pdf}
\end{itemize}
Hugo Schuchardt's 1909 article is the foundational description; an English translation appears in Gilbert (ed.), \emph{Pidgin and Creole Languages: Selected Essays by Hugo Schuchardt} (Cambridge University Press, 1980). No online copy was found. No dictionary was found.

\section{Religious and Community Resources}
None found. The language has no speaking community.

\section{Digital and Audio Resources}
No audio, video or software was found.

\section{Curriculum Development Context}
Because the attested material is sparse, a heritage curriculum would be a history and language-contact unit rather than a language course. Students could look at sample attested sentences, how the pidgin formed, and its Mediterranean trade setting, using the scholarly sources above. Real-language use in class would need reconstruction beyond what the sources support.
\end{document}